% Compile with XeLaTeX. This file is standalone; references are mapped in section 4 of the package.
\documentclass[UTF8,a4paper,11pt,fontset=fandol]{ctexart}
\usepackage[left=23mm,right=23mm,top=25mm,bottom=24mm,headheight=16pt]{geometry}
\usepackage{fontspec}
\setmainfont{lmroman10-regular.otf}[BoldFont=lmroman10-bold.otf,ItalicFont=lmroman10-italic.otf,BoldItalicFont=lmroman10-bolditalic.otf]
\setsansfont{lmsans10-regular.otf}[BoldFont=lmsans10-bold.otf,ItalicFont=lmsans10-oblique.otf,BoldItalicFont=lmsans10-boldoblique.otf]
\setmonofont{lmmono10-regular.otf}[Scale=0.84,ItalicFont=lmmono10-italic.otf,BoldFont=lmmonolt10-bold.otf]
\usepackage{amsmath,amssymb,mathtools,bm,esint,mathrsfs}
\usepackage{longtable,booktabs,array,calc,etoolbox}
\usepackage{graphicx}
\usepackage{fancyvrb}
\usepackage{xurl}
\usepackage{fancyhdr,lastpage}
\usepackage[unicode,hidelinks]{hyperref}
\hypersetup{pdftitle={来源、继承关系与推导定位},pdfauthor={Qiushi Engine research handoff},pdfsubject={Quantitative Einstein compatibility and mathematical physics}}
\setlength{\parindent}{0pt}
\setlength{\parskip}{0.45em}
\setlength{\emergencystretch}{3em}
\setlength{\tabcolsep}{5pt}
\renewcommand{\arraystretch}{1.22}
\AtBeginEnvironment{longtable}{\small}
\setcounter{secnumdepth}{0}
\setcounter{tocdepth}{2}
\ctexset{
 section={format=\Large\bfseries\raggedright,beforeskip=1.3em,afterskip=0.6em},
 subsection={format=\large\bfseries\raggedright,beforeskip=1em,afterskip=0.4em},
 subsubsection={format=\normalsize\bfseries\raggedright,beforeskip=0.8em,afterskip=0.3em}}
\pagestyle{fancy}
\fancyhf{}
\fancyhead[L]{\small\sffamily Qiushi Engine · 第一科学主线}
\fancyhead[R]{\small 2026-09-28}
\fancyfoot[C]{\small \thepage\ / \pageref{LastPage}}
\renewcommand{\headrulewidth}{0.3pt}
\renewcommand{\footrulewidth}{0pt}
\providecommand{\tightlist}{\setlength{\itemsep}{0pt}\setlength{\parskip}{0pt}}
\providecommand{\pandocbounded}[1]{#1}
\newcommand{\handoffpart}[1]{\clearpage\section{#1}}
\begin{document}
\begin{titlepage}
\thispagestyle{empty}
\vspace*{2cm}
{\large\sffamily QIUSHI ENGINE\par}
\vspace{1.3cm}
{\Huge\bfseries 来源、继承关系与推导定位\par}
\vspace{0.8cm}
{\Large 原稿快照与公开文献的可追溯入口\par}
\vspace{1.2cm}
{\large 跨会话研究综合与证明方案\par}
\vfill
{\normalsize 原稿状态基线：2026年9月27日，快照 \texttt{844e62f}\par}
\vspace{0.35cm}
{\normalsize 整理日期：2026年9月28日\par}
\vspace{1cm}
{\small 保留已取得的强结论，沿真实方程继续推进全局数学与物理问题。\par}
\end{titlepage}
\tableofcontents
\clearpage
\section{1. 内部研究依据}\label{ux5185ux90e8ux7814ux7a76ux4f9dux636e}

\textbf{{[}S0{]} 原稿包。}
\texttt{Yau26\_Line1\_Recent\_Research\_20260927.zip}，快照
\texttt{844e62f}。包含六组证明阅读卷及导读，MD与TeX是同一内容的两种格式，不是两套独立成果。本次整理直接读取了包内与
I40/I41、核心代价、首项及流相关的正文段落；本套文件不是对全部历史原稿的重新证明。

重要定位：

\begin{longtable}[]{@{}
  >{\raggedright\arraybackslash}p{(\columnwidth - 2\tabcolsep) * \real{0.5000}}
  >{\raggedright\arraybackslash}p{(\columnwidth - 2\tabcolsep) * \real{0.5000}}@{}}
\toprule\noalign{}
\begin{minipage}[b]{\linewidth}\raggedright
原稿包文件
\end{minipage} & \begin{minipage}[b]{\linewidth}\raggedright
与本套文件的关系
\end{minipage} \\
\midrule\noalign{}
\endhead
\bottomrule\noalign{}
\endlastfoot
\texttt{00\_READ\_FIRST.md} &
当前状态、旧状态被后稿替代的说明、I44/I45/I47/I49/X06边界 \\
\texttt{01\_linear\_gap\_and\_actual\_applications.md} & Source 1
§§4--5：I40；Source 2 Part I §§1--4：I41、共形约化、BF实际新窗口 \\
\texttt{02\_endpoint\_covering\_and\_sharp\_subfamilies.md} &
几何零端、覆盖估计、只添柄子族精确下确界与局部实际模型 \\
\texttt{03\_full\_cores\_signed\_extrema\_and\_flow.md} & Source 1
§1：完整核心与边界恒等式；Source 2
§§1、3--6：流、X06、Hessian与受限提取 \\
\texttt{04\_analytic\_foundations.md} &
实际场、二阶预算与前置解析工具 \\
\texttt{05\_actual\_topology\_and\_prior\_field\_equations.md} &
实际拓扑供给及先前场方程 \\
\texttt{06\_review\_records\_and\_source\_map.md} &
原路径、来源映射与历史核查记录 \\
\end{longtable}

\textbf{{[}S1{]} 独立审读报告。}
\texttt{Yau26\_Line1\_External\_Review\_20260927.md}，2026-09-27。主要使用§§二--五的核心定理、证明机制、三个开放连接和经典比较，以及§七的固定流形量词。

\textbf{{[}S2{]} 统一任务书。}
\texttt{Yau26\_Line1\_Research\_Directive\_20260927.md}，2026-09-27，接续
I40/I41 与
\texttt{e58b867}。主要使用§§3--7的弱积分零端、高指数模型、同尺度首项、标记传递及新范畴目标。

\textbf{{[}S3{]} 深化报告。}
\texttt{Qiushi\_Einstein\_研究主线深化与证明方案\_20260926.md}。用于经典继承表、复双曲边界恒等式、多实际场判据、固定
\(E_D\) 与慢模态路线。其§8曾将线性阶列为待证，已被{[}S0--S2{]}中的
I40/I41 更新；不得恢复旧状态。

\texttt{e58b867} 的完整源码不在本次可读取的 \texttt{844e62f}
原稿包中，因此高指数局部模型按{[}S2{]}中的摘要接续。最新
\texttt{LINE1\_CURRENT.md} 和
\texttt{20260927\_marked\_charge\_and\_first\_order\_flux.md}
应由拥有工作区的接收会话读取；本套文件不虚构它们的后续内容。

\section{2.
公开数学与物理文献}\label{ux516cux5f00ux6570ux5b66ux4e0eux7269ux7406ux6587ux732e}

以下标号与正文一致。外部文献用作经典背景、模型定义或最近邻比较；它们不是
I40/I41 新证明的替代证据。

\textbf{{[}R1{]} Claude LeBrun.} \emph{Ricci Curvature, Minimal Volumes,
and Seiberg--Witten Theory.} Inventiones Mathematicae 145 (2001),
279--316. arXiv:math/0003068. 重点比较：原论文的曲率能量、等号／极值和
Einstein 障碍，尤其原稿复用的 Theorem 4.3 拼接输入。
\href{https://arxiv.org/abs/math/0003068}{原文入口}

\textbf{{[}R2{]} D. Kotschick.} \emph{Monopole Classes and Einstein
Metrics.} International Mathematics Research Notices 2004(12), 593--609.
arXiv:math/0306022. 重点比较：special monopole classes 与 Theorem 12
的完整假设。 \href{https://arxiv.org/abs/math/0306022}{原文入口}

\textbf{{[}R3{]} Masashi Ishida and Claude LeBrun.} \emph{Curvature,
Connected Sums, and Seiberg--Witten Theory.} arXiv:math/0111228.
重点比较：BF实际类供给、Yamabe／曲率变分和 Einstein
障碍；不能只将其归为纯拓扑代数。
\href{https://arxiv.org/abs/math/0111228}{原文入口}

\textbf{{[}R4{]} Claude LeBrun.} \emph{Four-Manifolds, Curvature Bounds,
and Convex Geometry.} arXiv:math/0611450.
重点比较：单极子类凸包、双手性限制与边界情况；必须区分已知类集合与完整类集合。
\href{https://arxiv.org/abs/math/0611450}{原文入口}

\textbf{{[}R5{]} David M. J. Calderbank and Michael A. Singer.}
\emph{Einstein Metrics and Complex Singularities.} Inventiones
Mathematicae 156 (2004), 405--443. arXiv:math/0206229.
用途：I47锐利上界的标量平坦 Kähler ALE
模型；该经典构造不是本项目新增构造。
\href{https://arxiv.org/abs/math/0206229}{原文入口}

\textbf{{[}R6{]} M. F. Atiyah, N. J. Hitchin and I. M. Singer.}
\emph{Self-Duality in Four-Dimensional Riemannian Geometry.} Proceedings
of the Royal Society of London A 362 (1978), 425--461.
用途：四维手性曲率分解与诱导联络的经典结构。
\href{https://royalsocietypublishing.org/rspa/article/362/1711/425/14679/Self-duality-in-four-dimensional-Riemannian}{期刊入口}

\textbf{{[}R7{]} Edward Witten.} \emph{Monopoles and Four-Manifolds.}
Mathematical Research Letters 1 (1994), 769--796. arXiv:hep-th/9411102.
用途：SW方程与四维超对称规范理论的经典联系；并不提供本项目所设想的一般定量
BPS 间隙。 \href{https://arxiv.org/abs/hep-th/9411102}{原文入口}

\textbf{{[}R8{]} Claude LeBrun.} \emph{The Einstein--Maxwell Equations,
Extremal Kähler Metrics, and Seiberg--Witten Theory.} arXiv:0803.3734.
用途：真实 Einstein--Maxwell 变分问题及其与 SW
曲率估计的直接关系，是物理迁移必须比较的最近邻工作。
\href{https://arxiv.org/abs/0803.3734}{原文入口}

\textbf{{[}R9{]} K. S. Stelle.} \emph{Renormalization of
Higher-Derivative Quantum Gravity.} Physical Review D 16 (1977),
953--969.
用途：高导数／二次曲率引力的基础模型背景；不用于声称本项目解决量子一致性问题。
\href{https://doi.org/10.1103/PhysRevD.16.953}{期刊入口}

\textbf{{[}R10{]} Alberto Salvio.} \emph{Quadratic Gravity.} Frontiers
in Physics 6 (2018), 77. arXiv:1804.09944.
用途：二次曲率作用量、耦合与高导数问题的研究综述；与{[}R9{]}共同固定物理语境。
\href{https://arxiv.org/abs/1804.09944}{原文入口}

\textbf{{[}R11{]} Guido Festuccia and Nathan Seiberg.} \emph{Rigid
Supersymmetric Theories in Curved Superspace.} JHEP 06 (2011), 114.
arXiv:1105.0689.
用途：真实曲空间超对称依赖辅助场与完整变分，不能由``非Einstein''直接推出超对称破缺。
\href{https://arxiv.org/abs/1105.0689}{原文入口}

\textbf{{[}R12{]} Andrea Braides.} \emph{A Handbook of
\(\Gamma\)-Convergence.}
用途：变分极限的状态空间、下极限和恢复序列等正式要求。此文献为{[}S2{]}给出的定义性来源；本套文件没有宣称已证明当前几何问题的
\(\Gamma\) 极限。
\href{https://cvgmt.sns.it/media/doc/paper/57/Handbook.pdf}{作者稿入口}

\textbf{{[}R13{]} Masashi Ishida, Shinichiroh Matsuo and Nobuhiro
Nakamura.} \emph{Yamabe Invariants and the
\(\mathrm{Pin}^-(2)\)-monopole Equations.} Mathematical Research Letters
23(4) (2016). arXiv:1505.05948.
用途：扭曲单极子供给、Yamabe／Einstein障碍与相应流结论。
\href{https://arxiv.org/abs/1505.05948}{原文入口}

\textbf{{[}R14{]} American Physical Society.} \emph{About Physical
Review Letters}, acceptance
criteria，访问日期2026-09-28。用途：核对重要物理问题、关键方法和广泛物理意义的编辑标准，不预测录用。
\href{https://journals.aps.org/prl/about}{官方页面}

\textbf{{[}R15{]} American Physical Society.} \emph{About Physical
Review X}, acceptance
criteria，访问日期2026-09-28。用途：核对基础发现、关键跨学科联系和显著推进的编辑标准，不预测录用。
\href{https://journals.aps.org/prx/about}{官方页面}

\textbf{{[}R16{]} Maciej Dunajski, Jan Gutowski, Wafic Sabra and Paul
Tod.} \emph{Cosmological Einstein--Maxwell Instantons and Euclidean
Supersymmetry: Anti-Self-Dual Solutions.} Classical and Quantum Gravity
28 (2011), 025007. arXiv:1006.5149.
用途：一个具有明确方程和经典分类的四维欧氏超引力入口。
\href{https://arxiv.org/abs/1006.5149}{原文入口}

\textbf{{[}R17{]} Marco M. Caldarelli and Dietmar Klemm.} \emph{All
Supersymmetric Solutions of N=2, D=4 Gauged Supergravity.} JHEP 09
(2003), 019. arXiv:hep-th/0307022. 用途：Killing
spinor、完整场方程与精确分类的 Lorentzian 比较；不直接当作欧氏定理使用。
\href{https://arxiv.org/abs/hep-th/0307022}{原文入口}

\textbf{{[}R18{]} Mauro Carfora, Claudio Dappiaggi, Nicolò Drago and
Paolo Rinaldi.} \emph{Ricci Flow from the Renormalization of Nonlinear
Sigma Models in the Framework of Euclidean Algebraic Quantum Field
Theory.} arXiv:1809.07652. 用途：一圈 sigma-model 重整化流与 Ricci flow
的联系。 \href{https://arxiv.org/abs/1809.07652}{原文入口}

\section{3.
本套文件新增整理的演绎与研究建议}\label{ux672cux5957ux6587ux4ef6ux65b0ux589eux6574ux7406ux7684ux6f14ux7eceux4e0eux7814ux7a76ux5efaux8bae}

以下内容区别于原稿编号，便于接收会话追踪。

\begin{longtable}[]{@{}
  >{\raggedright\arraybackslash}p{(\columnwidth - 4\tabcolsep) * \real{0.3333}}
  >{\raggedright\arraybackslash}p{(\columnwidth - 4\tabcolsep) * \real{0.3333}}
  >{\raggedright\arraybackslash}p{(\columnwidth - 4\tabcolsep) * \real{0.3333}}@{}}
\toprule\noalign{}
\begin{minipage}[b]{\linewidth}\raggedright
内容
\end{minipage} & \begin{minipage}[b]{\linewidth}\raggedright
性质
\end{minipage} & \begin{minipage}[b]{\linewidth}\raggedright
定位
\end{minipage} \\
\midrule\noalign{}
\endhead
\bottomrule\noalign{}
\endlastfoot
\(\mathcal E\) 的显式公式与标量方差分解 &
直接代数推导，原项目已有相关关系 & 第一份文件式(5) \\
\(P+D_2\ge cAh\) 的明确常数形式 & 从 I41 与原稿恒等式直接推出 &
第二份文件式(M7) \\
正耦合二次曲率作用量间隙 & 本次补充的直接推论；物理查新仍需完成 &
第三份文件式(P7) \\
曲率手性块范数、诱导联络解释 & 经典结构的明确归一化与应用解释 &
第三份文件式(P1)--(P2) \\
Einstein--matter应力与Maxwell手性下界 &
条件于原几何定理及选定真实场方程的直接推论 & 第三份文件式(P11)--(P15) \\
一圈流去整体尺度的残差公式 & 直接投影计算，不是完整RG固定点结论 &
第三份文件式(P23) \\
缺陷有效理论、一般BPS间隙、全模态紧化 & 待证研究目标 &
第二份§9、第三份§§9--13 \\
\end{longtable}

简单恒等式可以精确复算，但它们不能替代 I40/I41
的解析证明，也不能自行认证后续物理目标的全球原创性。接收会话应保留这种依赖方向。

\section{4.
继续研究时的最小查新合同}\label{ux7ee7ux7eedux7814ux7a76ux65f6ux7684ux6700ux5c0fux67e5ux65b0ux5408ux540c}

对每个拟称为新结果的命题，记录：最近邻原文定理；完整输入；当前定理恢复了什么；严格新增的量词、对象或数值范围；一个真实实例；尚不能被该比较覆盖的文献范围。

若目标是解决一项物理猜想，应取得它的原始或公认现代准确陈述，再证明映射到当前框架。若做到了新的物理结论但没有解决那个猜想，应保留结果自身的价值，不把名称借来提高等级。
\end{document}
